Aby bolo možné tieto dáta efektívne analyzovať, bolo potrebné ich prekonvertovať pomocou basecallera Dorado do formátu BAM, ktorý je štandardným formátom pre ukladanie sekvenačných dát a umožňuje efektívne spracovanie a analýzu.
Po konverzii do formátu BAM bolo potrebné dáta zoskupiť a zarovnať, aby sa identifikovali podobné sekvencie a vytvorili konsenzuálne sekvencie, ktoré by mohli reprezentovať pôvodný DNA reťazec.
